%%%%%%%%%%%%%%%%%%%%%%%%%%%%%%%%%%%%%%%%%%%%%%%%%%%%%%%%%%%%%%%%%%%%%%
% Lớp 10 | Đại số tổ hợp | Hoán vị. Chỉnh hợp. Tổ hợp
% Dạng 0D8?2-9: Hoán vị lặp
%%%%%%%%%%%%%%%%%%%%%%%%%%%%%%%%%%%%%%%%%%%%%%%%%%%%%%%%%%%%%%%%%%%%%%


% nguồn: TuLuan_Lop10_Dot2 (nhập 14/11/2026)
\begin{bt}%[0D8C2-9]%[Dự án đề kiểm tra HKII NH22-23- Dương Phước Sang]%[THPT Lương Thế Vinh - Quảng Nam]%[ma:947a52e2]
		Từ các chữ số của tập $A=\big\{1;2;3;4;5;6;7\big\}$ có thể lập được bao nhiêu số tự nhiên có $8$ chữ số sao cho trong số có $8$ chữ số của số được lập ra mỗi chữ số của tập $A$ đều có mặt ít nhất một lần và không có hai chữ số chẵn nào đứng cạnh nhau.
		\loigiai{
			Tập $A$ có $7$ phần tử, để lập được số có $8$ chữ số mà các chữ số của $A$ đều xuất hiện ta phải dùng một chữ số $x$ nào đó trong $A$ $2$ lần.
			\begin{itemize}
				\item Trường hợp $1$: $x$ lẻ, tức là $x \in \big\{1;3;5;7\big\}$. Có $4$ cách chọn $x$. \\
				Để không có $2$ chữ số chẵn nào kề nhau ta làm như sau:\\
				$\bullet$ Trước hết, xếp $5$ chữ số $x$, $1$, $3$, $5$, $7$ thành hàng ngang, có $\dfrac{5!}{2!}$ cách xếp.\\
				\phantom{$\bullet$} Có $4$ vị trí nằm giữa mỗi cặp hai chữ số lẻ kề nhau được tạo ra.\\
				$\bullet$ Xếp $3$ chữ số $2$, $4$, $6$ vào $3$ trong $4$ vị trí nêu trên, có $\mathrm{A}_4^3$ cách xếp.\\
				Đến đây, dãy các chữ số đã xếp được hình thành nên một số thỏa mãn yêu cầu đề bài.\\
				Trường hợp $1$ ta lập được $4\cdot \left(\dfrac{5!}{2!}\cdot \mathrm{A}_4^3\right)=5760$ số như thế. 
				\item Trường hợp $2$: $x$ chẵn, tức là $x \in \big\{2;4;6\big\}$. Có $3$ cách chọn $x$. \\
				Để không có $2$ chữ số chẵn nào kề nhau ta chỉ cần xếp $4$ chữ số lẻ ($1$, $3$, $5$, $7$) và $4$ chữ số chẵn ($x$, $2$, $4$, $6$) xen kẽ nhau. Có $2 \cdot 4! \cdot \dfrac{4!}{2!}$ cách xếp chúng xen kẽ.\\
				Như vậy, trường hợp $2$ ta lập được $3\cdot \left(2 \cdot 4! \cdot \dfrac{4!}{2!}\right)=1728$ số thỏa mãn yêu cầu đề bài.
			\end{itemize}
			Vậy tổng cộng có $4\cdot \left(\dfrac{5!}{2!}\cdot \mathrm{A}_4^3\right)+3\cdot \left(2 \cdot 4! \cdot \dfrac{4!}{2!}\right)=5760+1728=7488$ số thỏa mãn yêu cầu đề bài.
		}
	\end{bt}

% nguồn: TuLuan_Lop10_Dot2 (nhập 14/11/2026)
\begin{bt}%[0D8H2-9]%[Dự án đề HKI-HKII K10 NH22-23 (Cả nước), Nhật Thiện]%[Số 1 - Sa Pa]%[ma:646afd16]
	Với các chữ số $0$, $1$, $2$, $3$, $4$, $5$, $6$. Lập được bao nhiêu số có mười chữ số mà trong mỗi số chữ số $3$ có mặt đúng $4$ lần, các chữ số khác mỗi chữ số có mặt đúng $1$ lần.
\loigiai{
Chọn $6$ trong $10$ vị trí xếp các chữ số $0$, $1$, $2$, $4$, $5$, $6$. Khi đó, ta có $A_{10}^6$ cách.\\
Xếp chữ số $3$ vào $4$ vị trí trống còn lại có $1$ cách.\\
Suy ra có $A_{10}^6\cdot 1=151\,200$ cách.\\
Xét trường hợp số $0$ đứng vị trí đầu tiên.\\
Chọn $5$ trong $9$ vị trí xếp các chữ số $1$, $2$, $4$, $5$, $6$. Khi đó, ta có $A_{9}^5$ cách.\\
Xếp chữ số $3$ vào $4$ vị trí trống còn lại có $1$ cách.\\
Suy ra có $A_9^5\cdot 1=15120$ cách.\\
Vậy có $151200-15120=136080$ số thỏa mãn bài toán.
}
\end{bt}
